% ------------------------------------------------------------
% Page 88
% ------------------------------------------------------------

Un peu plus loin, nous pouvons prendre sur la gauche, un chemin, maintenant goudronné, qui
nous amène au hameau de \textbf{Crouzet}, deux ou trois maisons, serrées les unes contre les autres,
sur le bord du ruisseau.

Le récit d’une assemblée qui se tint à la Saint Michel 29 Septembre 1697, dans la maison de
Cabanel, l’aîné, nous est parvenu. « Rey, \textit{(compagnon du prédicant Roman) fit une prière et
Roman prêcha. Sept ou huit hommes placés devant la porte, le fusil à la main, refusèrent de
laisser entrer deux jeunes filles de Pourcarès sous prétexte qu’elles dénonceraient les
assistants à Poujol (le traître de Pourcarès) ».}

\vspace{0.2cm}

En aval du Crouzet subsiste encore un petit moulin à farine et on trouve assez facilement les
entrées de galeries de mines de baryte, exploitées jusque vers les années 1950.

\vspace{0.2cm}

Un kilomètre seulement nous sépare de \textbf{Pourcarès} (nom qui signifie probablement la maison
des porcs, « porcos area » en latin), dominé par son fameux monolithe qui figurait dans les
livres de géographie du début du siècle dernier (Lou roun de la quille). Comme dans toute la
vallée, des assemblées, durement réprimées, se tinrent dans les environs, entre 1680 et au
moins 1702, comme au Villaret, dans la vallée voisine où les habitants furent condamnés à
payer 200 livres, après l’assemblée du 23 mars 1702, dans la maison tenue par Aiguebonne et
son gendre Martin.

En effet « \textit{le 4 Mai 1702, moi, Jacques Champagny, huissier royal au Vigan me suis porté à
cheval au lieu du Villaret au domicile de Pierre Gasse ( Gache ?), du dit lieu et lui ai intimé
et signifié la requête de M. l’Intendant...} »

Pourcarès avait connu une certaine animation tant que furent exploitées les mines de plomb-
argentifère jusqu’en 1918 environ. Mais aujourd’hui encore, la qualité et le goût de l’eau de sa
fontaine attirent les connaisseurs.

\vspace{0.2cm}

La vallée devient plus étroite mais très vite s’élargit pour accueillir \textbf{Connilergues}, ferme
isolée (le champ aux lapins (conillum ager), dont nous connaissons mieux l’histoire grâce aux
documents qui nous ont été remis par Françoise et Gilles Coupet-Hubac.

\vspace{0.2cm}

\textbf{Le domaine de Connilergues} est une dépendance de la baronnie de Roquedols au 14° siècle
et, plus précisément, « \textit{ancienne résidence du sire de Pourcarès} », dit une chronique du temps.
Sur la carte de Cassini les maisons de Connilergues sont surmontées d’une petite oriflamme
indiquant la noblesse des propriétaires.

\vspace{0.2cm}

Dans les années 1500, en effet, Pierre de Pagès se déclare « \textit{seigneur de Pourycrayres,
Roquedolz, le Villaret et Counilhargues} ». Beaucoup plus tard, vers 1675, Moïse et Samuel
Poujol du Crouzet - les prénoms bibliques de l’Ancien Testament montrent que la Réforme
est passée par là - sont « rentiers » (gérants) des mas de Connilergues, Campredon et Campis.
Ils perçoivent les rentes et les redevances, pour un certain « Abel de Maillan », autre titre du
Seigneur de Pourcayrès et de Roquedolz, et, probablement, de la majeure partie des terrés
de la vallée de la Brèze et du Béthuzon.
